\begin{abstract}

    \begin{center}
        \large{Использование генеративных нейронных сетей для трансфера изображений радиолокационного диапазона в оптический} \\
    \large\textit{Селин Тимур Александрович} \\[1 cm]

    Настоящая работа посвящена исследованию методов трансляции изображений радиолокационного диапазона в оптический с использованием генеративных нейронных сетей. Рассматривается задача перевода данных радиолокационной съёмки с синтезированной апертурой в изображения видимого спектра, что позволяет повысить интерпретируемость всепогодных спутниковых данных.

    В ходе выполнения работы проведён анализ нескольких общедоступных наборов парных изображений SAR-Оптика, основное внимание уделено датасету SEN12. Выполнено экспериментальное тестирование ряда генеративных архитектур: апробирована базовая модель Pix2Pix, а также специализированная архитектура CFRWD‑GAN, предназначенная для решения задачи SAR‑to‑Optical перевода. На основе полученных результатов предложены и реализованы доработки, направленные на повышение качества синтезируемых изображений.

    Полученные в работе результаты могут быть использованы в системах обработки данных дистанционного зондирования Земли для улучшения визуального анализа радиолокационной съёмки.
    
    Ключевые слова: генеративные нейронные сети, SAR-to-Optical, трансляция изображений, дистанционное зондирование Земли, GAN, SEN12.

    % \vfill

    % \textbf{Abstract} \\[1 cm]

    % Research and development of machine learning methods
    \end{center}

\end{abstract}
\newpage